\subsection{Quantum Corrections and the de Sitter Swampland Conjecture --- arXiv:1808.07498}

\textbf{Authors:} Keshav Dasgupta, Maxim Emelin, Evan McDonough, Radu Tatar

\paragraph{主な主張}
M理論の運動方程式と量子補正を調べ、考察したde Sitter解は一般に制御された四次元有効理論を欠く一方、穏やかな時間依存性が補正の階層を生みquasi-de Sitter解を許す可能性を示す\cite{Dasgupta:2018rtp}。

\paragraph{新規性}
ポテンシャルへの推測的条件だけでなく、on-shellの高次元運動方程式と量子補正の階層性からde Sitter予想を検討する。

\paragraph{位置付け}
時間依存背景における加速膨張と、有効理論の打ち切りが有効である条件との関係を調べる研究。

\paragraph{コメント（読書メモ）}
読書メモでは、(1) de Sitter予想と(2) 距離予想を組み合わせ、「(2)を満たし(1)を満たさない場合はlarge-field inflation、両方を満たす場合はquintessence」という対応として整理した（模型ごとの成立条件は要確認）。また、de Sitter予想に反する真空、あるいは漸近的なde Sitter状態が存在する場合に、四次元有効理論も構成できるのか、という二つの問いに注目した。本論文の設定では、時間非依存の場合に無限個の量子補正が同程度に寄与して有効理論の制御が失われる一方、有効な宇宙項の時間発展を許すとquasi-de Sitter構成の可能性が生じる、と理解した。
